%! Sine and Cosine Function Graphs — Unit 3, Lesson 3.4 — Homework
\documentclass[10pt]{article}
\usepackage{precalculus-article}
\usepackage{precalculus-boxes}
\newcommand{\TallMath}[1]{$\displaystyle #1\rule[-1.4em]{0pt}{3.2em}$}

\begin{document}

\pageheader{Unit 3, Lesson 3.4}{Homework}
\namedateperiod

\vspace{0.12in}

\begin{enumerate}[leftmargin=1.5em, itemsep=14pt]

  \item The function $f(\theta) = \sin\theta$ completes one full cycle on $[0, 2\pi]$.

        \begin{enumerate}[leftmargin=1.5em, itemsep=8pt]
            \item How many complete cycles does $f(\theta) = \sin\theta$ complete on
                  $[0, 6\pi]$? \blank{3cm}

            \item How many zeros does $f(\theta) = \sin\theta$ have on $[0, 4\pi]$?
                  (Count $\theta=0$ and $\theta=4\pi$.) \blank{3cm}

            \item On which interval(s) within $[0, 2\pi]$ is $\sin\theta > 0$?

                  \blank{8cm}
        \end{enumerate}

  \item \textbf{True or False?} The range of $f(\theta) = \cos\theta$ is $[0, 1]$.
        Circle your answer, then explain.

        \vspace{4pt}
        \textbf{TRUE} \quad / \quad \textbf{FALSE}

        \vspace{6pt}
        \writelines{2}

  \item Without computing, explain why $\sin(\theta + 2\pi) = \sin\theta$ for any angle
        $\theta$. Your explanation should reference the unit circle.

        \vspace{4pt}
        \writelines{3}

  \item \textbf{AP-Style Multiple Choice.} Which of the following intervals contains a
        zero of $f(\theta) = \cos\theta$?

        \vspace{4pt}
        \begin{enumerate}[label=(\Alph*), leftmargin=1.5em, itemsep=4pt]
            \item $\left(0,\;\dfrac{\pi}{4}\right)$
            \item $\left(\dfrac{\pi}{4},\;\dfrac{\pi}{2}\right)$
            \item $\left(\dfrac{\pi}{2},\;\pi\right)$
            \item $\left(\pi,\;\dfrac{3\pi}{2}\right)$
        \end{enumerate}

  \item The maximum value of $f(\theta) = \sin\theta$ is \blank{2cm} and it first occurs
        at $\theta = $ \blank{2cm}. The next time the function reaches that same maximum
        is at $\theta = $ \blank{2cm}.

  \item On the interval $[0, 2\pi]$, the function $f(\theta) = \cos\theta$ is
        \textbf{decreasing} on the interval \blank{5cm} and \textbf{increasing} on the
        interval \blank{5cm}.

\end{enumerate}

\vspace{0.1in}

\begin{extensionbox}
\textbf{Optional Extension:} On what interval(s) in $[0, 2\pi]$ is $f(\theta) = \cos\theta$
concave down (bending downward)? Relate your answer to the arc of the unit circle: where is
the unit-circle point moving in a way that causes the $x$-coordinate to decrease at an
increasing rate?

\vspace{6pt}
\writelines{4}
\end{extensionbox}

\vspace{0.1in}

\begin{blurbbox}[Preview: Lesson 3.5]{sky}
In Lesson 3.5 we will transform the basic graphs of $\sin\theta$ and $\cos\theta$ by
changing three parameters: \textbf{amplitude} (how tall the wave is), \textbf{period}
(how wide one cycle is), and \textbf{midline} (where the wave is centered vertically).
These transformations let us model real-world phenomena such as tides, temperature cycles,
and sound waves.
\end{blurbbox}

\end{document}
